\chapter{Image không cần daemon: Jib và registry cục bộ}
\label{ch:jib}

\section{Vì sao không cứ \code{docker build} trong CI?}

Các Dockerfile vẫn chạy tốt. Ba lý do khiến pipeline build image theo
cách khác:

\begin{enumerate}
  \item \textbf{CI không cần Docker daemon.} \code{docker build} đòi hỏi
  một daemon đang chạy; cấp cho Jenkins một daemon nghĩa là
  docker-in-docker hoặc mount socket Docker của máy chủ --- tức là quyền
  root trên máy chủ, trao cho mọi build. Jib build và push image bằng
  Java thuần: không daemon, không đặc quyền.
  \item \textbf{Phân layer khớp với cách code thay đổi.} Dockerfile sao
  chép một fat jar: một layer khổng lồ, push lại toàn bộ mỗi lần build.
  Jib tách ứng dụng thành các layer theo mức độ biến động --- dependency
  / resource / class --- nên sửa một dòng code chỉ push lại layer class
  nhỏ. Trên hai nhân CPU và Wi-Fi, đó là khác biệt giữa vài giây và vài
  phút cho mỗi service.
  \item \textbf{Sẵn sàng cho OpenShift một cách tình cờ.} Image do Jib
  tạo chạy bằng người dùng không phải root với quyền file thuộc group 0
  --- đúng điều mô hình bảo mật UID tùy ý của OpenShift đòi hỏi. Tính
  khả chuyển miễn phí.
\end{enumerate}

\section{Cấu hình, có chú giải}

Jib được cấu hình một lần trong \code{pluginManagement} của
\code{pom.xml} cha; mỗi module chạy được tham gia bằng một tham chiếu
\code{<plugin>} trơn, còn \code{common} \emph{rút lui} bằng
\code{jib.skip=true}:

\begin{lstlisting}[language=xml]
<plugin>
  <groupId>com.google.cloud.tools</groupId>
  <artifactId>jib-maven-plugin</artifactId>
  <version>3.4.4</version>
  <configuration>
    <from><image>eclipse-temurin:25-jre-alpine</image></from>  (*@\co{1}@*)
    <to><image>${docker.registry}/${project.artifactId}:${project.version}</image></to> (*@\co{2}@*)
    <allowInsecureRegistries>true</allowInsecureRegistries>    (*@\co{3}@*)
    <container>
      <mainClass>${start-class}</mainClass>                    (*@\co{4}@*)
      <creationTime>USE_CURRENT_TIMESTAMP</creationTime>       (*@\co{5}@*)
    </container>
  </configuration>
</plugin>
\end{lstlisting}

\begin{description}
  \item[\co{1}] Cùng image gốc với các Dockerfile --- hành vi lúc chạy
  không đổi, chỉ khác đường đóng gói.
  \item[\co{2}] Phân giải thành \code{localhost:5000/course-service:1.0.0-SNAPSHOT}.
  Tên được viết \emph{từ góc nhìn của cụm} ---
  Mục~\ref{sec:whylocalhost}.
  \item[\co{3}] Registry cục bộ dùng HTTP thường; Jib từ chối push không
  an toàn trừ khi được bảo. Mọi công cụ trong chuỗi này (Jib, containerd)
  đều chọn an-toàn-mặc-định như nhau và đều cần cùng một lời cho phép
  tường minh.
  \item[\co{4}] Đặt tường minh vì bộ quét bytecode đi kèm Jib chưa đọc
  được file class của Java~25 --- một ví dụ thật về việc Java quá mới
  làm khổ bộ công cụ.
  \item[\co{5}] Jib mặc định dùng mốc thời gian epoch-0 để build tái lập
  được; ở đây ta đánh đổi điều đó để \code{docker images} hiển thị thời
  gian build thật.
\end{description}

\section{Registry, và vì sao lại là \code{localhost:5000}}
\label{sec:whylocalhost}

\begin{lstlisting}[language=bash]
docker run -d --restart=unless-stopped --name registry \
  -p 5000:5000 -v registry-data:/var/lib/registry registry:2
\end{lstlisting}

Sự thật không hiển nhiên buộc phải có kiến trúc này: \textbf{k3s không
nhìn thấy image của Docker.} k3s chạy containerd riêng với kho image
riêng; \code{docker build} trên máy chủ đặt image vào nơi k3s không bao
giờ ngó tới. Registry là giao diện trung thực giữa bên build và cụm ---
thứ mà pipeline đằng nào cũng cần.

Tên image không chỉ là cái tên --- nó là \emph{pull từ đâu}:
\code{localhost:5000/course-service} nghĩa là ``pull từ
localhost:5000.'' Vì bên pull chính là node k3s, và Jenkins (bên push)
dùng chung network namespace với máy chủ (Chương~22), cùng một chuỗi
này hợp lệ cho cả hai --- và \code{localhost} không bao giờ trôi như IP
LAN hay hostname có thể. Hai điểm cấu hình làm nó chạy được:

\begin{lstlisting}[language=yaml]
# /etc/rancher/k3s/registries.yaml -- containerd assumes HTTPS otherwise
mirrors:
  "localhost:5000":
    endpoint: ["http://localhost:5000"]
\end{lstlisting}

\begin{pitfall}[http: server gave HTTP response to HTTPS client]
Lỗi đặc trưng của chương này. Cả Jib (push) lẫn containerd (pull) đều
mặc định TLS; \code{registry:2} chỉ nói HTTP thường. Nếu một pod kẹt ở
\code{ImagePullBackOff} với dòng chữ này trong
\code{kubectl describe pod}, thì mục trong \code{registries.yaml} đang
thiếu --- và lưu ý containerd \emph{chỉ đọc nó lúc khởi động}:
\code{systemctl restart k3s} sau khi sửa.
\end{pitfall}

\begin{decision}[registry cục bộ vs.\ GHCR]
GitHub Container Registry sẽ cho phép pull image từ bất cứ đâu và thêm
quét lỗ hổng --- với cái giá là đẩy hàng trăm megabyte lên Wi-Fi gia
đình mỗi lần build. Registry trong LAN giữ vòng lặp build ở tốc độ ổ
đĩa. Quyết định này đảo chiều ngay khi có máy thứ hai cần pull image;
khi đó trong manifest không gì thay đổi ngoài tiền tố image.
\end{decision}

\section{Tag, digest, và cái bẫy SNAPSHOT}

\emph{Tag} (\code{:1.0.0-SNAPSHOT}, \code{:a1b2c3d}) là con trỏ có thể
dời; \emph{digest} (\code{@sha256:...}) là mã băm nội dung bất biến.
Pipeline push mỗi build với hai nhãn: tag SNAPSHOT (mà manifest tham
chiếu) và tag git-SHA (mà các lần deploy ghim vào --- Chương~23).
Kubernetes cache image theo tag; một SNAPSHOT push lại sẽ bị node lặng
lẽ bỏ qua --- đó là lý do mọi Deployment đều đặt
\code{imagePullPolicy: Always}: kiểm tra lại registry mỗi lần khởi
động; nếu digest khớp, phép kiểm tra chỉ tốn một lần tải manifest.

\begin{quiz}
\begin{enumerate}
  \item Nêu ba vấn đề CI mà Jib giải quyết so với build dựa trên daemon,
  và vấn đề nào là chuyện bảo mật.
  \item Vì sao tên image phải viết từ góc nhìn của \emph{cụm}, và điều
  gì hỏng nếu Jenkins và k3s hiểu \code{localhost:5000} khác nhau?
  \item Một đồng nghiệp đề xuất \code{docker save | k3s ctr images import}
  thay cho registry. Cách đó tốn gì mỗi lần deploy, và stage nào của
  pipeline tự nhiên ngừng hoạt động?
  \item Giải thích cách một \code{:1.0.0-SNAPSHOT} push lại cộng với
  pull policy mặc định tạo ra tình huống ``tôi đã deploy mà chẳng có gì
  đổi,'' và cả hai cách sửa.
\end{enumerate}
\end{quiz}
